% Options for packages loaded elsewhere
\PassOptionsToPackage{unicode}{hyperref}
\PassOptionsToPackage{hyphens}{url}
\documentclass[
  conference]{IEEEtran}
\usepackage{xcolor}
\usepackage{amsmath,amssymb}
\setcounter{secnumdepth}{5}
\usepackage{iftex}
\ifPDFTeX
  \usepackage[T1]{fontenc}
  \usepackage[utf8]{inputenc}
  \usepackage{textcomp} % provide euro and other symbols
\else % if luatex or xetex
  \usepackage{unicode-math} % this also loads fontspec
  \defaultfontfeatures{Scale=MatchLowercase}
  \defaultfontfeatures[\rmfamily]{Ligatures=TeX,Scale=1}
\fi
\usepackage{lmodern}
\ifPDFTeX\else
  % xetex/luatex font selection
\fi
% Use upquote if available, for straight quotes in verbatim environments
\IfFileExists{upquote.sty}{\usepackage{upquote}}{}
\IfFileExists{microtype.sty}{% use microtype if available
  \usepackage[]{microtype}
  \UseMicrotypeSet[protrusion]{basicmath} % disable protrusion for tt fonts
}{}
\makeatletter
\@ifundefined{KOMAClassName}{% if non-KOMA class
  \IfFileExists{parskip.sty}{%
    \usepackage{parskip}
  }{% else
    \setlength{\parindent}{0pt}
    \setlength{\parskip}{6pt plus 2pt minus 1pt}}
}{% if KOMA class
  \KOMAoptions{parskip=half}}
\makeatother
\setlength{\emergencystretch}{3em} % prevent overfull lines
\providecommand{\tightlist}{%
  \setlength{\itemsep}{0pt}\setlength{\parskip}{0pt}}
\usepackage{bookmark}
\IfFileExists{xurl.sty}{\usepackage{xurl}}{} % add URL line breaks if available
\urlstyle{same}
\hypersetup{
  hidelinks,
  pdfcreator={LaTeX via pandoc}}

\author{}
\date{}

\begin{document}

\section{ASE 2026 --- Neural-driven
Computing}\label{ase-2026-neural-driven-computing}

Bernhard Guillon

\subsection{Introduction}\label{introduction}

Modern large language models and AI toolchains typically run inside
heavy software stacks (OS + browser + runtime). Many applications will
increasingly embed neural components at many layers of the stack. This
project explores whether we can design a far more lightweight computing
environment that can directly host perceptron-based neural models as
first-class computational elements.

We propose to define a minimal instruction set (ISA) with a small set of
neural/tensor primitives, implement an emulator for that ISA, and run
small trained networks directly in the emulator. The effort covers ISA
design, assembler/loader, model-to-ROM/memory serialization, and
verification tooling (including automated black‑box tests).

\subsection{Vision}\label{vision}

Create a compact, demonstrable platform where:

A tiny neural model (e.g., a perceptron or small MLP) can be loaded and
executed by a minimal runtime / emulator.

The ISA is RISC-like (RISC-V inspired) and extended with a few
neural-tensor primitives (matmul, activation, tensor ops).

I/O is minimal: character input (keyboard-like device) and a simple
framebuffer output. The pipeline is: Keyboard input → neural model
inference → Framebuffer pixel output

The initial demonstration task: train a model that receives short
character input and produces a corresponding framebuffer output (e.g.,
render text or glyphs). The emulator should be capable of loading the
model into memory and running inference using the extended ISA.

\subsubsection{Use of AI}\label{use-of-ai}

Use LLM to help with research on the ISA, create AI training data. Also
write the emulator for the ISA a bootloader and an assembler. We might
also need a ``transpiler'' to transform the model to something which we
can load. Or write our own model creator.

\end{document}
